% ==============================================================================
% Law (Legal Scholarship) Paper Template
% Structure: IRAC (Issue / Rule / Application / Conclusion) and Law-Review
% style with heavy footnoted citations.
%
% Citation: Bluebook (US) / OSCOLA (UK) / AGLC (AU). Footnote-heavy; BibTeX
% automation is partial. This skeleton uses standard \footnote{} with
% manually formatted citations + plain.bst for any bibliography entries.
%
% Compile:  pdflatex main -> bibtex main -> pdflatex main -> pdflatex main
% ==============================================================================
\documentclass[11pt,a4paper]{article}

% ---- Page geometry ----
\usepackage[a4paper,top=25mm,bottom=25mm,left=30mm,right=25mm]{geometry}

% ---- Fonts & encoding ----
\usepackage[T1]{fontenc}
\usepackage[utf8]{inputenc}
\usepackage{mathptmx}
\usepackage{helvet}
\usepackage{courier}

% ---- Footnote styling: law reviews use footnotes heavily ----
\usepackage[bottom,hang,flushmargin]{footmisc}  % hang footnotes, no indent
\renewcommand{\footnoterule}{\vspace*{0.3em}\hrule width 0.4\textwidth\vspace*{0.3em}}
\setlength{\footnotesep}{0.8em}

% ---- Cross-references ----
\usepackage{cleveref}             % \cref{} for smart cross-references

% ---- Math, figures, tables ----
\usepackage{amsmath,amssymb}
\usepackage{graphicx}
\usepackage{booktabs}
\usepackage{array}
\usepackage{xcolor}

% ---- Citations: minimal BibTeX support; most cites go in \footnote{} ----
\usepackage[numbers,sort&compress]{natbib}
\bibliographystyle{plain}         % law reviews rarely use BibTeX for citations;
% use biblatex with style=oscola or style=british-legal if you want automation.

% ---- Hyperref last ----
% \usepackage{hyperref}
% \hypersetup{colorlinks=true,linkcolor=blue,citecolor=blue,urlcolor=blue}

% ---- Italic case names (legal convention) ----
\newcommand{\case}[1]{\textit{#1}}  % e.g. \case{Brown v. Board of Education}

\title{Title of the Legal Submission: \\ A Doctrinal and Policy Analysis}
\author{First Author\textsuperscript{*}\\
\textsuperscript{*}School of Law, University of Y, City, Country; e-mail: first.author@univ.edu}
\date{}

\begin{document}
\maketitle

% ==============================================================================
% Abstract  (~150-200 words; state the legal issue, thesis, and contribution)
% ==============================================================================
\begin{abstract}
% TODO: ~150-200 words. Sentence 1: the legal problem. Sentence 2-3: the
% gap in the doctrine or scholarship. Sentence 4-5: the thesis. Sentence 6:
% the proposed solution and its implications.
\end{abstract}

\noindent\textbf{Keywords:} % TODO: 4-6 keywords.
keyword one, keyword two, keyword three

% ==============================================================================
% 1. Introduction
% ==============================================================================
\section{Introduction}
\label{sec:intro}
% TODO: open with the legal problem and its societal stakes; state the
% thesis up front (law reviews favour front-loaded theses); preview the
% roadmap of the article. Cite authorities in footnotes, e.g.
% \footnote{See \case{Brown v. Board of Education}, 347 U.S. 483 (1954).}

% ==============================================================================
% 2. IRAC Analysis  (the doctrinal core)
% ==============================================================================
\section{Issue}
\label{sec:issue}
% TODO issue: frame the legal question as a "Whether..." statement.
% Example: \textit{Whether [party]'s conduct constitutes [tort/crime]
% under [statute] when [facts].}

\section{Rule}
\label{sec:rule}
% TODO rule: state the controlling statutes, restatements, and leading
% cases. Use \footnote{} for each citation in Bluebook form.
% Example: Under the Restatement (Second) of Torts \S 281 (1965), a
% plaintiff must prove duty, breach, causation, and damages.
% \footnote{\textit{Restatement (Second) of Torts} \S 281 (Am. Law Inst. 1965).}

\subsection{Statutory framework}
% TODO: quote and analyse the relevant statute(s).

\subsection{Case law}
% TODO: trace the doctrinal evolution through leading cases; distinguish
% majority / concurrence / dissent; note circuit splits.

\section{Application}
\label{sec:application}
% TODO application: apply the rule to the facts step by step. Use analogical
% reasoning to compare the facts at bar with those of leading cases; note
% policy considerations where the doctrine is indeterminate.

\subsection{Analogical reasoning from precedent}
% TODO: for each analogous case, summarise facts + holding + reasoning, then
% explain why the present facts are similar or distinguishable.

\subsection{Counter-arguments}
% TODO counter-argument: address the strongest counter-interpretation of
% the rule; explain why the proposed reading is superior.

\section{Conclusion}
\label{sec:conclusion}
% TODO conclusion: state the outcome of the IRAC analysis; flag any
% remaining ambiguities or circuit splits.

% ==============================================================================
% 3. Law-Review Style Sections  (for full-length law review articles)
% ==============================================================================
% The IRAC analysis above is the doctrinal core. The sections below extend
% it into a full law-review article with background, proposal, and
% implications. For a short case note, comment out sections 3-7.

\section{Background}
\label{sec:background}
% TODO: trace the historical development of the doctrine; situate the issue
% within the broader statutory and constitutional landscape.

\section{The Problem}
\label{sec:problem}
% TODO: identify the doctrinal gap, contradiction, or injustice that this
% article addresses; cite criticism in the scholarly literature.

\section{Proposed Solution}
\label{sec:proposal}
% TODO: articulate a legislative, judicial, or administrative reform;
% specify the textual changes to statutes, regulations, or doctrine.

\section{Implications}
\label{sec:implications}
% TODO: assess how the proposed reform would affect adjacent areas of law
% and the institutional roles of courts, legislatures, and agencies.

% ==============================================================================
% Final Conclusion  (for full law review articles)
% ==============================================================================
\section{Final Conclusion}
\label{sec:final}
% TODO: restate the thesis and the proposed solution in one paragraph.

% ==============================================================================
% Declarations
% ==============================================================================
\section*{Declarations}

\noindent\textbf{Funding:} % TODO: grant numbers and funding bodies.

\noindent\textbf{Conflicts of interest:} % TODO: declare or state none.

% ==============================================================================
% Acknowledgements
% ==============================================================================
\section*{Acknowledgements}
% TODO: funders, colleagues, librarians, research assistants.

% ==============================================================================
% References  —  supplementary (most citations are in footnotes)
% Compile:  pdflatex main -> bibtex main -> pdflatex main -> pdflatex main
% ==============================================================================
\bibliography{references}
% Law reviews rarely use a numbered bibliography; uncomment below only if the
% target venue requires a "References" list in addition to footnotes.
% \begin{thebibliography}{99}
% \bibitem{ref_example} Author, F. \emph{Title of the Book}. Publisher, Year.
% \end{thebibliography}

\end{document}
